\chapter{Conclusion}

The course followed one thread: the Bellman equations (\cref{ch:mdp}),
solved exactly when the model is known (\cref{ch:dp}), from samples when it
is not (\cref{ch:td}), with function approximation when the states are too
many (\cref{ch:vfa,ch:dqn}); then the direct optimisation of the policy
(\cref{ch:pg,ch:ppo}), extended to continuous actions (\cref{ch:continuous});
and planning with a model (\cref{ch:model-based}). The last chapters
removed, one at a time, the assumptions of the standard setting: rewards
replaced by demonstrations, interaction replaced by a fixed dataset
(\cref{ch:imitation}), exploration made deliberate (\cref{ch:exploration}),
several learning agents (\cref{ch:multiagent}), and rewards learnt from
human preferences (\cref{ch:rlhf}).

A few ideas came back in almost every chapter: bootstrapping, and the
contraction that makes it converge; the bias--variance trade-off between
Monte Carlo returns and bootstrapped targets, controlled by $n$-step returns
and $\lambda$; the danger of maximising over estimates (the maximisation
bias of DQN, the actor exploiting the critic in DDPG, extrapolation in
offline RL, reward hacking in RLHF), countered by pessimism or by keeping
close to a reference; and the value of keeping the policy close to the one
that produced the data (trust regions, KL penalties, behaviour constraints).

On Pacman, these ideas took the agent from $79\%$ to about $98\%$ of games
won: PPO with feature planes, a mirror-symmetric network and GAE, then a
look-ahead search that knows nothing of the ghosts' code, and expert
iteration to teach the search's decisions to the network. Game \code{v2}
(a bigger maze, three ghosts that come back after being eaten, a shorter
power) opens the next round, with a new difficulty: as discussed in
\cref{rem:pomdp}, it is only partially observable from one image. The
directions left open in \cref{ch:model-based} (MCTS with a value network,
an AlphaZero-style loop) and the learning ghosts of \cref{ch:multiagent}
are natural next projects.

\paragraph{Further reading.} Sutton and Barto, \emph{Reinforcement
Learning: An Introduction} (2nd edition, 2018) for the foundations;
Szepesv\'ari, \emph{Algorithms for Reinforcement Learning} (2010) and
Bertsekas, \emph{Reinforcement Learning and Optimal Control} (2019) for the
theory; Lattimore and Szepesv\'ari, \emph{Bandit Algorithms} (2020) for
bandits; Agarwal, Jiang, Kakade and Sun, \emph{Reinforcement Learning:
Theory and Algorithms} (online notes) for sample complexity; and the
original papers cited in each chapter. For Monte Carlo tree search in
two-player games, see the Hex implementation and the interactive
visualiser cited in \cref{rem:mcts-two-players}.
